%%%%%%%%%%%%%%%%%%%%%%%%%%%%%%%%%%%%%%%%%%%%%%%%%%%%%%%%%%%%%
%% Examtask 1: Step-up converter with output filter %%
%%%%%%%%%%%%%%%%%%%%%%%%%%%%%%%%%%%%%%%%%%%%%%%%%%%%%%%%%%%%%

\task{Step-down converter}

%%%%%%%%%%%%%%%%%%%%%%%%%%%%%%%%%%%%%%%%%%%%%%
\taskGerman{Tiefsetzsteller}

A drive system is powered by DC-voltage.
The speed, position, and current sensors in the drive system require stable lower voltage.
To transform the voltage down efficiently, a buck converter is used.

\begin{germanblock}
	Ein Antriebssystem wird mit Gleichstrom versorgt.
	Die Drehzahl-, Positions- und Stromsensoren im Antriebssystem benötigen eine niederigere, stabile Versorgungsspannung.
	Um die Spannung effizient herunterzutranformieren, wird ein Tiefsetzsteller verwendet.
\end{germanblock}

% figure figStepDownConverterOutputFilter 
\input{fig/Task01/figStepDownConverterOutputFilter}

\begin{table}[ht]
	\centering  % Zentriert die Tabelle
	\begin{tabular}{llll}
		\toprule
		\multicolumn{2}{l}{\textbf{General parameters:}} & \multicolumn{2}{l}{\textbf{IGBT:}}                                                                                               \\
		Input voltage:                                   & $U_{\mathrm{1}} = \SI{60}{\volt}$          & Collector-emitter voltage:          & $u_{\mathrm{on},\mathrm{CE}} = \SI{2}{\volt}$ \\
		Output voltage:                                  & $U_2 = \SI{5.1}{\volt}$                    &                                     &                                               \\
		Output power:                                    & $\SI{2.5}{\watt} \le p \le \SI{15}{\watt}$ & \multicolumn{2}{l}{\textbf{Diode:}}                                                 \\
		Switching frequency:                             & $f_\mathrm{s} = \SI{100}{\kilo\hertz}$     & Forward voltage:                    & $u_\mathrm{fw} = \SI{0.8}{\volt}$             \\
		\bottomrule
	\end{tabular}
	\caption{Parameters of the step-down converter.}  % Beschriftung der Tabelle
	\label{table:ex01_Parameters of the circuit}
\end{table}

\subtask{What duty cycle must be set for the given input-to-output voltage ratio, assuming ideal components?}{1}

\subtaskGerman{Welches Tastverhältnis muss für das gewünschte Spannungsübersetzungsverähltnis eingestellt werden, unter der Annahme idealer Bauteile?}

% Solution of subtask
\begin{solutionblock}
	The duty cycle corresponds to
	\begin{equation*}
		D=\frac{U_\mathrm{2}}{U_\mathrm{1}}=\frac{\SI{60}{\volt}}{\SI{5.1}{\volt}}=0.085.
	\end{equation*}
\end{solutionblock}

\subtask{What duty cycle is required when the voltage drops across the transistor and the diode are taken into account?}{3}

\subtaskGerman{Welches Tastverhältnis ist erforderlich, wenn die Spannungsabfälle über dem Transistor und der Diode berücksichtigt werden?}

% Solution of subtask
\begin{solutionblock}
	In steady state, the inductor current balance is to apply:
	\begin{equation*}
		\Delta i_\mathrm{L,T_\mathrm{on}} \cdot T_\mathrm{on}=\Delta i_\mathrm{L,T_\mathrm{off}} \cdot T_\mathrm{off}.
	\end{equation*}
	If the transistor conducts $\Delta i_\mathrm{L,T_\mathrm{on}}$ results in:

	\begin{equation*}
		\Delta i_\mathrm{L,T_\mathrm{on}}=\frac{(U_\mathrm{1}-U_\mathrm{T}-U_\mathrm{2}) \cdot T_\mathrm{on} }{L}.
	\end{equation*}

	If the transistor is blocking $\Delta i_\mathrm{L,T_\mathrm{off}}$ is calculated by:
	\begin{equation*}
		\Delta i_\mathrm{L,T_\mathrm{off}}=\frac{(U_\mathrm{2}+u_\mathrm{fw}) \cdot T_\mathrm{off}}{L}.
	\end{equation*}

	By substituting the currents within the inductor current balance equation, we obtain
	\begin{equation*}
		\frac{(U_\mathrm{1}-U_\mathrm{T}-U_\mathrm{2}) \cdot T_\mathrm{on}}{L}=\frac{(U_\mathrm{2}+u_\mathrm{fw}) \cdot T_\mathrm{off}}{L}.
	\end{equation*}

	Replacing $T_\mathrm{on}$ and $T_\mathrm{off}$ by the switching period yields to
	\begin{equation*}
		\frac{(U_\mathrm{1}-U_\mathrm{T}-U_\mathrm{2}) \cdot T_\mathrm{s} \cdot D}{L} =
		\frac{(U_\mathrm{2}+u_\mathrm{fw}) \cdot (1-D) \cdot T_\mathrm{s}}{L}.
	\end{equation*}

	After simplifying by canceling $T_\mathrm{s}$ and $L$, we obtain
	\begin{equation*}
		(U_\mathrm{1}-U_\mathrm{T}-U_\mathrm{2}) \cdot D =
		(U_\mathrm{2}+u_\mathrm{fw}) \cdot (1-D).
	\end{equation*}
	Solving this equation with respect to $D$ yields
	\begin{equation*}
		D = \frac{U_\mathrm{2}+u_\mathrm{fw}}{U_\mathrm{1}-U_\mathrm{T}+u_\mathrm{fw}} =
		\frac{\SI{5.1}{\volt}+\SI{0.8}{\volt}}{\SI{60}{\volt} - \SI{2}{\volt}+\SI{0.8}{\volt}}  = 0.1.
	\end{equation*}

\end{solutionblock}

\subtask{The sensors have different power consumption depending on the operating point, i.e., $P_2 \in [\SI{2.5}{\watt}, \SI{15}{\watt}]$ is given.
	Which load operating point is the relevant one to design the inductance value $L$ such that discontinuous conduction mode (DCM) is always avoided?
	Calculate the minimal required inductances $L$ for this case.}{2}

\subtaskGerman{Die Sensoren haben je nach Betriebspunkt einen unterschiedlichen Leistungsbedarf, d.h. es gilt $P_2 \in [\SI{2.5}{\watt}, \SI{15}{\watt}]$. Welcher Lastbetriebspunkt ist für die Auslegung der Induktivität $L$ relevant, um den Betrieb im Lückbetrieb zu vermeiden? Berechnen Sie die minimal erforderliche Induktivität $L$ für diesen Fall.}

% Solution of subtask
\begin{solutionblock}
	The BCM mode occurs when the minimum average inductor current equals the half of the inductor current ripple.
	This leads to the operation point with minimal power consumption, which corresponds to the minimal current.
	The minimal average current $\overline{i}_\mathrm{1}$ is expressed by
	\begin{equation*}
		\overline{i}_\mathrm{2,min} = \frac{P_\mathrm{2}}{U_\mathrm{2}} = \frac{\SI{2.5}{\watt}} {\SI{5.1}{\volt}} = \SI{0.49}{\ampere}.
	\end{equation*}
	In case of boundary conduction mode following equation is to apply:
	\begin{equation*}
		\Delta i_\mathrm{L}= 2 \cdot \overline{i}_\mathrm{1,min} = 2 \cdot \SI{0.49}{\ampere} = \SI{0.98}{\ampere}.
	\end{equation*}
	The ripple current $\Delta i_\mathrm{L}$ yields
	\begin{equation*}
		\Delta i_\mathrm{L}=\frac{U_\mathrm{L,on} \cdot T_\mathrm{on}}{L}
		=\frac{(U_\mathrm{1} - U_\mathrm{T} - U_\mathrm{2}) \cdot D}{L \cdot f_\mathrm{s}}.
	\end{equation*}
	Solving this equation with respect to $L$ yields
	\begin{equation*}
		L= \frac{(U_\mathrm{1} -U_\mathrm{T} - U_\mathrm{2}) \cdot D}{\Delta i_\mathrm{L} \cdot f_\mathrm{s}}
		= \frac{(\SI{60}{\volt} - \SI{2}{\volt} - \SI{5.1}{\volt}) \cdot 0.1} {\SI{0.98}{\ampere} \cdot \SI{100}{\kilo\hertz}}
		= \SI{54}{\micro\henry}.
	\end{equation*}
\end{solutionblock}

\subtask{In the event of a malfunction, the speed and position sensors fail. As a result, the current $I_\mathrm{2,fault}$ drops to a value of $\SI{0.25}{\ampere}$. Which conduction mode will the converter enter, and what are the resulting consequences?
	Calculate the output voltage for this fault condition. The voltage losses off the diode and transistor shall not be considered}{2}
\subtaskGerman{ In einem Störfall kommt es zum Ausfall der Drehzahl- und Positionssensoren. Dadurch fällt der Strom $I_\mathrm{2}$ auf einen Wert von $\SI{0.25}{\ampere}$.
	Wie hoch ist die Ausgangsspannung für diesen Störfall? Die Spannungsabfälle über den Transistor und die Diode sollen nicht berücksichtigt werden.}
% Solution of subtask
\begin{solutionblock}
	The output current is less than the minimal current in BCM. This leads to DCM, which causes an overvoltage.
	The output voltage is to be calculated by the equation for DCM:
	%\begin{equation*}
	\begin{align*}
		U_\mathrm{2} & = \frac{D^\mathrm{2} \cdot T_\mathrm{s} \cdot U_\mathrm{1}^\mathrm{2}} {D^\mathrm{2} \cdot T_\mathrm{s} \cdot U_\mathrm{1}^\mathrm{2}
		+ 2 \cdot L \cdot  \overline{i}_\mathrm{2,fault}}                                                                                                    \\
		             & = \frac{D^\mathrm{2} \cdot U_\mathrm{1}^\mathrm{2}} {D^\mathrm{2} \cdot U_\mathrm{1}^\mathrm{2}
		+ 2 \cdot L \cdot  \overline{i}_\mathrm{2,fault} \cdot f_\mathrm{s}}                                                                                 \\
		             & = \frac{0.085^\mathrm{2} \cdot \SI{60}{\volt}^\mathrm{2}} {0.085^\mathrm{2} \cdot \SI{60}{\volt}^\mathrm{2}
		+ 2 \cdot \SI{54}{\micro\henry} \cdot  \SI{0.25}{\ampere} \cdot \SI{100}{\kilo\hertz}}                                                               \\
		             & = \SI{8.3}{\volt}.
	\end{align*}
	%\end{equation*}
\end{solutionblock}
